\begin{conclusion}
\label{chapitre:conclusion}

\newcommand{\conclusionheading}[1]{%
  \par\medskip
  \noindent\textbf{#1}\par\smallskip
}

\conclusionheading{Synthèse}
Ce mémoire avait pour objectif de concevoir un système reproductible de détection de boiterie
probable à partir de données capteurs bovines, tout en conservant une logique interprétable et
scientifiquement défendable. Cet objectif est atteint.

Le système proposé combine un détecteur d'anomalies non supervisé (Isolation Forest) avec des
règles métier explicites (persistance temporelle, cohérence multi-signaux, espacement des
notifications). Cette architecture hybride transforme des signaux capteurs bruts
et bruités en épisodes de boiterie probable, interprétables et actionnables pour le suivi
d'élevage.

\conclusionheading{Rappel des résultats principaux}
La validation du système a été conduite à deux niveaux complémentaires:
\begin{itemize}
    \item \textit{Validation manuelle} (v2.1): les deux événements détectables injectés sont
    récupérés avec un recouvrement strict (IoU $\geq$ 0.20), tandis que l'événement
    volontairement non détectable reste silencieux. L'impact sur le troupeau est localisé à
    la seule vache cible (+2~notifications), les 27~autres animaux restant strictement
    stables;
    \item \textit{Validation robuste}: sur 9\,504~exécutions couvrant 336~configurations,
    le système détecte 100\% des cas forts en détection brute
    (\texttt{strong\_detect\_any\_mean}=1.000), atteint un recouvrement strict moyen de
    0.667 sur ces cas forts, maintient une fuite nulle sur les scénarios non détectables,
    et contrôle les fausses notifications à 0.105 par vache et par jour. L'accord de décision entre les campagnes v5.1 et v5.2 atteint
    91.7\%, confirmant la stabilité de la calibration;
    \item \textit{Ablation élargie}: l'étude sur 90~cas par variante montre que chaque
    composante du pipeline est nécessaire. L'IF seul détecte tout mais sans précision
    temporelle; les règles seules génèrent $3.2\times$ plus de fausses alertes. LOF + règles
    obtient le meilleur compromis empirique sur ce benchmark, ouvrant une piste
    d'amélioration crédible. Les écarts entre variantes sont confirmés par des tests
    statistiques appariés (McNemar exact, Wilcoxon signed-rank);
    \item \textit{Comparaison à une baseline et calibration}: la comparaison à la
    baseline pédométrique historique d'Alsaaod (2012) s'appuie sur des
    intervalles de confiance à 95\,\% par bootstrap (B=20\,000), et les scores
    bruts du détecteur sont recalibrés par régression isotonique avec IC95 par bootstrap
    (B=2\,000), atteignant Brier~=~0.087 [IC95\,\%: 0.083; 0.091] et
    ECE~=~0.015 [IC95\,\%: 0.010; 0.021];
    \item \textit{Interprétabilité}: une analyse SHAP (TreeExplainer sur Isolation Forest)
    et des courbes ROC/précision-rappel complètent les règles métier en documentant
    l'importance relative des variables dérivées utilisées par le détecteur.
\end{itemize}

Le benchmark de performance confirme une scalabilité linéaire ($\approx$0.6~seconde par
vache), compatible avec un déploiement en temps raisonnable sur des troupeaux de taille
courante. Une matrice RACI documente l'intégration opérationnelle du système dans le
workflow de ferme (éleveur, vétérinaire, technicien, système).

\conclusionheading{Contribution}
La contribution principale ne réside pas dans l'invention d'un nouvel algorithme théorique,
mais dans l'assemblage rigoureux d'un système appliqué et sa validation défendable. Le projet
délivre:
\begin{enumerate}
    \item un pipeline complet et modulaire, depuis l'ingestion des données capteurs
    jusqu'aux notifications;
    \item un protocole de validation en deux niveaux, documentant explicitement les forces et
    les limites du système;
    \item une étude d'ablation honnête, qui ne masque pas les alternatives crédibles, et
    étayée par des tests appariés (McNemar exact, Wilcoxon signed-rank);
    \item une comparaison à une baseline pédométrique historique (Alsaaod 2012) avec
    intervalles de confiance à 95\,\% par bootstrap, et une calibration isotonique des scores
    assortie des métriques Brier et ECE;
    \item une couche d'interprétabilité (SHAP TreeExplainer, courbes ROC/précision-rappel)
    complémentaire aux règles métier;
    \item une interface interactive (dashboard Streamlit) rendant les alertes directement
    exploitables par l'utilisateur final (visualisation par vache, classement du troupeau,
    carte de chaleur temporelle);
    \item une matrice RACI documentant l'intégration opérationnelle du système dans le
    workflow de ferme;
    \item une traçabilité expérimentale complète (SHA-256), garantissant que les résultats
    présentés sont reproductibles et auditables.
\end{enumerate}

\conclusionheading{Limites}
La principale limite du système est l'absence de validation clinique synchrone:
les performances rapportées reposent sur des injections synthétiques, pas sur des cas
de boiterie confirmés par un vétérinaire. Une validation exploratoire sur le corpus
complémentaire IceTag 2019 (30~vaches de la même ferme, 12~avec scores SLS) a confirmé
la portabilité technique du pipeline et une cohérence partielle avec les antécédents
cliniques, sans toutefois permettre de conclusion clinique en raison du décalage de
8 à 9~mois entre les évaluations et les données capteurs
(Section~\ref{sec:validation_mcgill}). Cette limite
est mise en perspective par la profondeur du protocole de validation interne: 9\,504~exécutions sur
336~configurations, ablation à 4~variantes, stabilité inter-seeds, et traçabilité
cryptographique des résultats. La fragilité sur les cas borderline, la dépendance aux
seuils calibrés sur un seul corpus, et la généralisation à d'autres fermes et capteurs
restent des axes d'amélioration. Ces limites sont partagées par la majorité des travaux
publiés dans le domaine.

\conclusionheading{Disponibilité}
Le code source du pipeline, les scripts de validation et les paramètres gelés
sont versionnés dans un dépôt Git privé. La version de référence du mémoire
correspond au tag \texttt{memoire-v2}, accompagné des manifestes d'intégrité
SHA-256 des artefacts expérimentaux. Le code et les artefacts non confidentiels
sont disponibles sur demande auprès de l'auteur. Les données brutes capteurs
ne peuvent être diffusées pour des raisons de confidentialité, mais les profils
d'injection, les paramètres de chaque campagne et la procédure de reproduction
(Annexe~\ref{annexe:validation}) permettent de reproduire intégralement la
démarche sur de nouvelles données.

\conclusionheading{Perspectives}
Plusieurs axes de travail futur sont identifiés:
\begin{enumerate}
  \item \textit{Validation multi-fermes}: étendre la validation à de nouveaux corpus
  provenant de fermes et de systèmes capteurs différents, afin de tester la
  généralisabilité des seuils et de l'architecture;
  \item \textit{Labels cliniques}: intégrer des labels cliniques mieux synchronisés lorsque
  des collaborations vétérinaires le permettent, afin de calculer des métriques de
  performance au sens clinique strict;
  \item \textit{Ensemble hybride IF--LOF}: la validation exploratoire (protocole v4,
  288~exécutions) révèle des profils complémentaires entre les deux détecteurs: IF plus précis
  temporellement sur les cas forts, LOF plus sensible sur les cas borderline
  (Chapitre~\ref{chapitre:discussion}). Un ensemble hybride combinant ces forces constitue
  une piste d'amélioration crédible, à valider sur les protocoles v5.1 et v5.2;
  \item \textit{Enrichissement des signaux}: explorer l'ajout de variables physiologiques
  (production laitière, température, poids) et environnementales (météo, saison) pour
  améliorer la spécificité de la détection;
  \item \textit{Déploiement opérationnel}: évaluer le système dans un contexte d'utilisation
  réelle, avec des retours d'éleveurs et de vétérinaires sur l'utilité et l'acceptabilité
  des alertes produites.
\end{enumerate}

\end{conclusion}
